% TOP_PROBLEM: {"id": 1243, "title": "Magic Square of Squares", "category_id": 1, "category": {"id": 1, "name": "number_theory", "display_name": "Number Theory", "description": "Properties of integers, prime numbers, Diophantine equations.", "slug": "number-theory", "order_index": 1, "created_at": "2026-07-31T15:26:25.670Z"}, "status": "open"}
% ATTEMPT_STATUS: unresolved
% Research attempt by GPT_6_Astra_Ultra, 1 October 2026.
% Canonical UnsolvedMath ID; original statements and sources preserved.
% FIRST_POSTED: 2026-10-01
% Imported from: unsolvedmath_additional_500.tex; UnsolvedMath ID 1243
% !TeX program = lualatex
% Compile twice: lualatex 1243.tex
% Self-contained: no external bibliography, images, or \input files.
% Numeric section labels and visible numbers are original UnsolvedMath IDs.
\documentclass[11pt,a4paper]{article}
\usepackage[margin=24mm,headheight=14pt]{geometry}
\usepackage{fontspec}
\setmainfont{Latin Modern Roman}
\setsansfont{Latin Modern Sans}
\setmonofont{Latin Modern Mono}
\usepackage{amsmath,amssymb,mathtools}
\usepackage{unicode-math}
\setmathfont{Latin Modern Math}
\usepackage{microtype}
\usepackage{enumitem,xurl,needspace}
\usepackage[dvipsnames]{xcolor}
\usepackage{hyperref}
\hypersetup{unicode=true,colorlinks=true,linkcolor=MidnightBlue,urlcolor=MidnightBlue,
 pdftitle={UnsolvedMath Problem 1243},pdfauthor={GPT 6 Astra Ultra},bookmarksnumbered=true}
\usepackage{bookmark,tocloft,titlesec,fancyhdr}
\setlength{\cftsecnumwidth}{4.2em}
\cftsetpnumwidth{2.5em}
\cftsetrmarg{4em}
\titleformat{\section}{\Large\bfseries\raggedright}{\thesection}{0.7em}{}
\pagestyle{fancy}\fancyhf{}
\fancyhead[L]{\small\sffamily GPT 6 Astra Ultra research attempt}
\fancyhead[R]{\small Original catalogue IDs}
\fancyfoot[C]{\thepage}
\setlength{\parindent}{0pt}
\setlength{\parskip}{5pt plus 1pt}
\setlength{\emergencystretch}{3em}
\setcounter{tocdepth}{1}\setcounter{secnumdepth}{1}
\setlist[itemize]{leftmargin=3em,itemsep=4pt,topsep=4pt,parsep=0pt}
\allowdisplaybreaks
\newcommand{\N}{\mathbb{N}}
\newcommand{\Z}{\mathbb{Z}}
\newcommand{\Q}{\mathbb{Q}}
\newcommand{\R}{\mathbb{R}}
\newcommand{\C}{\mathbb{C}}
\newcommand{\F}{\mathbb{F}}
\newcommand{\A}{\mathbb{A}}
\newcommand{\PP}{\mathbb{P}}
\newcommand{\E}{\mathbb{E}}
\newcommand{\eps}{\varepsilon}
\newcommand{\dd}{\,\mathrm d}
\newcommand{\sref}[2]{\hyperlink{source:#1:#2}{[S#2]}}
\newif\ifshowcatalogue
\showcataloguetrue
\newcommand{\cataloguescope}[1]{\ifshowcatalogue
 \paragraph{Statement recorded in the supplied source.}
 #1\fi}
\begin{document}
% UnsolvedMath ID 1243; source catalogue code NT-036

\setcounter{section}{1242}

\section{A Three-by-Three Magic Square of Distinct Squares}\label{1243}

{\small\sffamily UnsolvedMath ID 1243 · Number Theory}

\subsection{Definitions and mathematical statement}

\paragraph{Shared notation.}Unless a section says otherwise, $\N=\{1,2,\ldots\}$,
$\Z$ is the ring of integers, $\Q,\R,\C$ are the rational, real, and complex
fields, $[N]=\{1,\ldots,\lfloor N\rfloor\}$, and $\log$ is the natural
logarithm. For real functions of a parameter tending to infinity,
$f\ll g$ and $f=O(g)$ mean $|f|\leq Cg$ eventually for some fixed $C$;
$f\gg g$ reverses this comparison; $f=o(g)$ means $f/g\to0$;
$f\sim g$ means $f/g\to1$; and $f\asymp g$ means both order bounds.
A subscript on $O$ or $\ll$ specifies allowed constant dependence.
The expression $n^{o(1)}$ in an upper bound means growth smaller than
$n^\epsilon$ for each fixed $\epsilon>0$. Local definitions govern any
exception, finite-field convention, or use of the same symbol for another
object. An asymptotic claim is not automatically an effective algorithm.



Seek integers $b_{ij}$ for $1\leq i,j\leq3$ such that the nine numbers $a_{ij}=b_{ij}^2$ are pairwise distinct and all three row sums, all three column sums and both main diagonal sums equal one common number. The wording perfect squares'' is used literally, with zero permitted unless a source explicitly imposes positivity. \sref{1243}{1} \sref{1243}{2}

\paragraph{Statement recorded in the supplied source.}
Does there exist a 3\ensuremath{×}3 magic square composed entirely of distinct perfect squares? \sref{1243}{1}

\subsection{Short English statement}

Can nine different square numbers be arranged in a three-by-three array so that every row, column and main diagonal has the same sum?

\subsection{Research attempt}

\paragraph{Scope and source check.}
All nine entries must be distinct integer squares, and all eight lines
must have the same sum. The literal statement permits zero, so it is
not silently excluded. Bremner's primary paper [E1] distinguishes
repeated-square magic examples and examples satisfying fewer than all
eight line sums; neither meets this target. The route here is a complete
linear normal form followed by exact square-pair constraints.

\paragraph{Every three-by-three magic array has three parameters.}
Let the central entry be $e$ and the common sum be $S$. Add the middle
row, middle column, and the two diagonals. This counts each noncentral
entry once and the centre four times, so the total is $3S+3e$.
It also equals $4S$, giving $S=3e$. Each opposite pair then sums to
$2e$. Taking the upper left and upper right entries to be $e+x$ and
$e+y$, the remaining line equations force the array
\[
 \begin{pmatrix}
 e+x&e-x-y&e+y\\
 e-x+y&e&e+x-y\\
 e-y&e+x+y&e-x
 \end{pmatrix}.
\]
Conversely direct addition verifies all eight sums are $3e$ for any
$e,x,y$. Thus no magic constraint is lost. The square problem asks for
$e=z^2$ and for every value in
\[
 \{e,\ e\pm x,\ e\pm y,\ e\pm(x+y),\ e\pm(x-y)\}
\]
to be a square, with this set having cardinality nine. Merely finding
four opposite square pairs with a common midpoint is not enough: their
offsets must also have the additive form $x,y,x+y,x-y$.

\paragraph{Zero is impossible even under the literal statement.}
If one entry is zero, its opposite entry is $2e=2z^2$. For $z\ne0$
this cannot be an integer square: writing $w^2=2z^2$ makes the
two-adic valuation on the left even and on the right odd. If $z=0$,
then $S=0$ and every nonnegative row entry must be zero, violating
distinctness. If zero were the centre, the same last argument applies.
Hence every actual solution automatically has nine positive entries.
No positivity assumption has been added to the catalogue statement.

\paragraph{Primitive congruences force a common residue.}
Divide all square roots by their common gcd to obtain a primitive array.
Its centre root $z$ cannot be even. If it were, $2e\equiv0\pmod4$;
the only way two square residues modulo four sum to zero is $0+0$.
Every opposite pair would then have even roots, contradicting
primitivity. With $z$ odd, $2e\equiv2\pmod8$. The square residues
$0,1,4$ have sum two modulo eight only as $1+1$, so every root is odd.

Similarly $3$ cannot divide $z$. If it did, opposite pairs would sum
to zero modulo three; two square residues can do that only as $0+0$,
making every root divisible by three. Thus $e\equiv1\pmod3$, and
every opposite pair has sum two modulo three, forcing both entries to
be one. Combining these facts, all nine entries in a primitive solution
are one modulo $24$. In particular
\[
 e\equiv1\pmod{24},\qquad x\equiv y\equiv0\pmod{24}.
\]
These congruences are necessary. They do not force the offsets to vanish
as integers; any inference from a residue class to equality would be
invalid.

\paragraph{Opposite pairs are rational points on a circle.}
After dividing an opposite square pair by $z^2$, its roots $u,v$ satisfy
$u^2+v^2=2$. The line through $(1,1)$ with rational slope $t$ gives
the second intersection
\[
 u=\frac{t^2-2t-1}{1+t^2},\qquad
 v=\frac{1-2t-t^2}{1+t^2}.
\]
Substitution verifies $u^2+v^2=2$; the tangent choice returns the base
point, and a vertical line handles its omitted direction. Signs of
the roots do not affect their square entries. Thus it is easy to
parametrize one opposite pair rationally. What remains is to obtain
four pairs whose normalized offsets close under both addition and
subtraction, while keeping the nine square values distinct. Four
independent rational circle parameters do not automatically satisfy
those two additive compatibility conditions.

\paragraph{An exact finite algorithm with a complete bound.}
For a fixed positive centre root $z$, define the finite signed set
\[
 D_z=\{t^2-z^2:0\le t\le\lfloor\sqrt{2z^2}\rfloor,
                       \ 2z^2-t^2\text{ is a square}\}.
\]
It is symmetric under negation and contains zero. The normal form gives
the exact criterion: a square exists with this centre if and only if
there are $x,y\in D_z$ such that $x+y,x-y\in D_z$ and
$\{0,\pm x,\pm y,\pm(x+y),\pm(x-y)\}$ has nine elements.
Indeed all required entries are nonnegative squares by the definition
of $D_z$, and the preceding normal form gives the magic sums.
Conversely every such array supplies these parameters. This criterion
checks distinctness and every line, not only a selected subset.

The script enumerates these sets for every $1\le z\le300$ and tests
every signed pair $(x,y)$. It finds no solution. For orientation,
$|D_1|=1$, $|D_5|=3$, $|D_{25}|=5$, and $|D_{65}|=9$.
Already having nine available offsets at $z=65$ is insufficient; the
additive closure and distinctness tests still fail. The computation
covers every array with centre root at most $300$, even if another
root exceeds $300$, since the enumeration allows roots up to
$\lfloor\sqrt2\,z\rfloor$. Files are \texttt{checks.py} and
\texttt{results.json} in
\path{_Local/Erdos_problems/UnsolvedMath/attempt_audit_2026_10_01/number_theory/}.

\paragraph{Remaining gap and outcome.}
The normal form is complete, zero has been ruled out without changing
the statement, and an exact finite-centre criterion has been executed.
No theorem prevents additive configurations of the required type in
$D_z$ for all larger $z$, and no finite upper bound on $z$ is proved.
\textbf{Outcome: UNRESOLVED.} A rational parametrization of one pair,
or a repeated-entry magic square, does not meet the nine-distinct-square
target.

\subsection{Sources}

{\small

\begin{itemize}

\item[\hypertarget{source:1243:1}{S1}] ULAM.AI, \emph{UnsolvedMath}, supplied \texttt{problems.json}, record 1243 (NT-036), statement and background. The embedded literature-review date is 2026-08-17; that is the archive’s review date, not a fresh certification. Dataset landing page: \url{https://huggingface.co/datasets/ulamai/UnsolvedMath}.

\item[\hypertarget{source:1243:2}{S2}] A. Bremner, On squares of squares, Acta Arithmetica 88 (1999), 289-297. \url{https://eudml.org/doc/207247}. Bibliographic information imported from the supplied record.

\item[\hypertarget{source:1243:3}{S3}] A. Bremner, On squares of squares II, Acta Arithmetica 99 (2001), 289-308. \url{https://doi.org/10.4064/aa99-3-6}. Bibliographic information imported from the supplied record.

\item[\hypertarget{source:1243:4}{S4}] Open Problem Garden, Magic square of squares. \url{https://www.openproblemgarden.org/op/magic_square_of_squares}. Bibliographic information imported from the supplied record.


\item[E1] Andrew Bremner, \emph{On squares of squares},
Acta Arithmetica 88 (1999), 289--297.
\url{https://matwbn.icm.edu.pl/ksiazki/aa/aa88/aa8837.pdf}.
Primary paper opening and formulation inspected 1 October 2026.
The deductions and finite-centre enumeration here are provided directly.

\end{itemize}

}

\label{end:1243}
\end{document}
